\documentclass[11pt]{article}

% ---------------------------------------------------------------------------
%  L5 — Convex polyplets.
%  Category L (entirely machine-written); see paper/README.md and
%  docs/publication-split.md.
%
%  N3 COLLISION.  Gouyou-Beauchamps & Leroux, FPSAC 2004, §2.3, have the block
%  decomposition of \label{lem:blocks} and the mirror equality that now lives in
%  L7, for convex polyominoes on the honeycomb lattice.
%  docs/publication-split.md requires that citation in THREE places, marked
%  "attribution --- k of n".  Since the L5/L7 split separated the two results,
%  the placements are: HERE rem:attr1 (at lem:blocks) and the paragraph opening
%  \label{sec:related}; in L7 rem:attr2 (at prop:mirror) and its own
%  \label{sec:related} paragraph.  Four placements, none droppable, and the
%  total must not fall below three.  docs/publication-split.md still describes
%  the pre-split arrangement and needs jasonp's sign-off to match.
%  (Labels, not compiled numbers, on purpose: numbering drifts under revision,
%  labels do not.)
%
%  Source material, all in this repository:
%    results/hv-growth-sandwich.md      Props 6-11, Lemmas 1-4, Conjecture 8
%    results/convex-polyplets.md        the exclusions, mu, the PSLQ boxes
%    docs/proofs/convex-mirage.md       the original mirage finding
%    results/middle-kingdom-phase3.md   the grid, Corollary 4, Proposition 5
%    results/middle-kingdom.md          the campaign index
%    results/mk-dir4-perimeter.md       the dir4 perimeter series
% ---------------------------------------------------------------------------

% ---------------------------------------------------------------------------
%  paper/shared/preamble.tex — packages and macros common to every manuscript
%  in this directory.  Factored out of paper/polyplets-report.tex and
%  paper/technical-report.tex, which between them had two copies of the same
%  twenty lines.
%
%  technical-report.tex does NOT \input this file and is not to be edited to do
%  so: it is jasonp's prose and is read-only to the machine (paper/README.md,
%  "Who may edit what").  The duplication there is deliberate.
%
%  Assumes \documentclass[11pt]{article} has already been issued.
% ---------------------------------------------------------------------------
\usepackage[margin=1in]{geometry}
\usepackage{amsmath,amssymb}
\usepackage{amsthm}
\usepackage{booktabs}
\usepackage{graphicx}
\usepackage{numprint}
\usepackage{microtype}
\usepackage[table]{xcolor}
\usepackage[hidelinks]{hyperref}
\usepackage{flafter}   % a float never appears before the text that introduces it
\usepackage{float}     % [H] pins tables/figures exactly where written
\usepackage[font=small,labelfont=bf,labelsep=period]{caption}
\setlength{\intextsep}{16pt plus 3pt minus 3pt}

\npthousandsep{,}
\npfourdigitsep
\newcommand{\num}[1]{\numprint{#1}}
\newcommand{\g}{\cellcolor{gray!15}}

% Theorem environments.  One counter shared across kinds, so that a reference
% like "Theorem 4" is unambiguous in a paper that also has lemmas and
% propositions.
\theoremstyle{plain}
\newtheorem{theorem}{Theorem}
\newtheorem{proposition}[theorem]{Proposition}
\newtheorem{lemma}[theorem]{Lemma}
\newtheorem{corollary}[theorem]{Corollary}
\theoremstyle{definition}
\newtheorem{definition}[theorem]{Definition}
\newtheorem{example}[theorem]{Example}
\theoremstyle{remark}
\newtheorem{remark}[theorem]{Remark}
\newtheorem{conjecture}[theorem]{Conjecture}
\newtheorem{openproblem}[theorem]{Open Problem}

% Repo-path citations.  Every one of these must resolve in the tree that ships
% with the paper; `make gate-citations` checks the markdown, and
% paper/verify_claims.py checks the manuscripts.
\newcommand{\repo}[1]{\texttt{#1}}

% OEIS links.
\newcommand{\oeis}[1]{\href{https://oeis.org/#1}{#1}}

% Growth constants, so a rename is one edit.
\newcommand{\lam}{\lambda}
\newcommand{\muH}{\mu_H}

% ---------------------------------------------------------------------------
%  paper/shared/disclosure.tex — the two authorship disclosure blocks.
%
%  docs/publication-split.md §1 fixes the rule these implement:
%
%    "Under no circumstances will there be any ambiguity about whether a paper
%     is wholly my work (none), merely my prose (some) or entirely LLM-authored
%     (some)."   — jasonp, 2026-08-07
%
%  There are exactly two blocks and they live in exactly one file, so that a
%  paper cannot quietly soften its own disclosure.  The fixed sentences take no
%  arguments.  The ONE thing a paper supplies is the per-result verification
%  ledger, because that is the part that differs honestly between papers — and
%  §1 requires it to be per result, not in aggregate.
%
%  Do not add a \Pdisclosure or \Ldisclosure variant with an optional softening
%  argument.  If a paper does not fit one of these two blocks, the paper is
%  wrong, not the block.
% ---------------------------------------------------------------------------

\newsavebox{\disclosurebox}

\newenvironment{disclosureframe}
  {\begin{center}\begin{lrbox}{\disclosurebox}\begin{minipage}{0.93\textwidth}\small}
  {\end{minipage}\end{lrbox}\fbox{\usebox{\disclosurebox}}\end{center}}

% --- P: jasonp's prose -------------------------------------------------------
% Byline: Jason Parker, alone.  Argument: the per-result ledger of what the
% machine did.
\newcommand{\Pdisclosure}[1]{%
\begin{disclosureframe}
\textbf{Authorship.}\enspace Every sentence of this paper was written by its
named author. He has read every proof it states and believes each one, and he
is responsible for the correctness of what it claims.

\smallskip
\textbf{Machine contribution.}\enspace #1
\end{disclosureframe}}

% --- L: entirely machine-written --------------------------------------------
% Argument: the per-result verification ledger — for each result, what warrants
% it (Lean kernel, exact certificate, or labelled measurement) and how much of
% it a human has checked.
\newcommand{\Ldisclosure}[1]{%
\begin{disclosureframe}
\textbf{Authorship.}\enspace This document was written end to end by a large
language model. That includes its mathematics, its prose, its tables and its
\LaTeX{}. No sentence of it was written by a human. The mathematics it reports
was likewise derived by machine.

\smallskip
\textbf{Human role.}\enspace The person who directed this work chose what to
compute, chose what to publish, and is responsible for the decision to release
this document. Except where the ledger below says otherwise, that person has
not personally checked the arguments here, and this disclosure is not to be
read as an endorsement of them.

\smallskip
\textbf{What warrants each result.}\enspace #1
\end{disclosureframe}}

% --- Draft banner ------------------------------------------------------------
% For an L-paper whose verification ledger is not yet settled.  Loud on purpose:
% an L-paper without a truthful ledger must not be mistaken for a finished one.
%
% \firstdraftbanner is the standard wording, which five papers previously
% carried character-for-character in their own preambles.  It lives here for the
% same reason the disclosure blocks do: identical text in seven places drifts,
% and a paper must not be able to soften it by editing its own copy.  The
% optional argument is for a gate specific to one paper, appended verbatim; a
% paper needing different wording (L6, compute-gated) calls \draftbanner direct.
\newcommand{\firstdraftbanner}[1][]{%
\draftbanner{This is a first draft. The verification ledger below records what
warrants each result \emph{mechanically}; the column that records what a human
has read is empty, deliberately, because nobody has read it yet. Do not
circulate until that column says something true.#1}}

\newcommand{\draftbanner}[1]{%
\begin{center}
\fcolorbox{red!60!black}{yellow!15}{%
  \begin{minipage}{0.93\textwidth}\small
  \textbf{DRAFT — NOT FOR CIRCULATION.}\enspace #1
  \end{minipage}}
\end{center}}


\newcommand{\Z}{\mathbb{Z}}
\newcommand{\Q}{\mathbb{Q}}
\newcommand{\R}{\mathbb{R}}
\newcommand{\HV}{\mathrm{HV}}

\title{\textbf{Convex polyplets}\\[4pt]
\large Intractable by area, algebraic by perimeter, and one growth
constant for every class in between}
\author{[byline: jasonp's call --- \repo{docs/publication-split.md} \S1]}
\date{August 2026}

\begin{document}
\maketitle

\firstdraftbanner[ One further gate is specific to this paper: the attribution
to Gouyou-Beauchamps and Leroux (\S\ref{sec:related}) must survive every
revision, in both of the places it appears here and in both of the places it
appears in companion paper L7.]

\Ldisclosure{%
\emph{Proposition~\ref{prop:squeeze} (the squeeze) and Lemmas~\ref{lem:blocks},
\ref{lem:stacks}, \ref{lem:supermul}} --- proved below, with no analytic input:
an injection, Fekete's lemma, and an elementary partition bound.
Lemma~\ref{lem:supermul}'s column join and its inverse are formalized in Lean~4
against mathlib (\texttt{Stair.join\_valid}, \texttt{Stair.cut\_join},
\texttt{Stair.join\_injOn}), with axiom footprint \texttt{[propext,
Quot.sound]} --- not even \texttt{Classical.choice} --- and the counting layer
and Fekete step are there too (\texttt{Stair.M\_supermul},
\texttt{Stair.M\_tendsto}, \texttt{Stair.M\_le\_mu\_pow}), all guarded. So
Lemma~\ref{lem:supermul} is a Lean theorem.
\emph{Proposition~\ref{prop:squeeze} itself is NOT formalized}: the geometric
HV-convex layer (Lemma~\ref{lem:blocks}) and the squeeze are paper proofs.
\emph{The exclusions of \S\ref{sec:exclusions}} --- rigorous in the direction
claimed: full column rank modulo a prime is a proof of full rank over $\Q$.
Four positive and negative control arms, and a fifth on an unrelated series,
are pinned by a gate.
\emph{$\mu$ to 199 digits, $\theta = 0$, the amplitudes} --- labelled
measurement, with the trusted digit count itself measured rather than asserted.
\emph{$q_c$, $\mu$ and $A$ to 44--45 digits (\S\ref{sec:kernel})} --- certified,
by exact rational interval arithmetic with no floating point in the chain; those
digits agree with the extrapolated ones. The certificate bounds the constants;
it does not by itself prove the sharp asymptotic.
\emph{Algebraicity of the area moments at every level (\S\ref{sec:moments})} ---
proved, by an explicit recursive kernel construction.
\emph{The area limit law} --- \textbf{published before us}
(\cite{richard2007,entingGuttmann1989}), rederived here unawares, and claimed
only in the narrow form \S\ref{sec:moments} states.
\emph{The zeros of $K$ and the non-D-finiteness route} --- \textbf{a lead}, firm
numerically, with two named gaps (Open Problem~\ref{op:kzeros}).
\emph{The PSLQ verdicts} --- labelled measurement, stated as falsifiable box
exclusions and never as ``$\mu$ is transcendental''.
\emph{The subdominant exponential and the amplitude ratio} --- moved to
companion paper L7, with their own ledger.
\emph{Novelty} --- N1 and N2 clean; \textbf{N3 is a real collision}, cited here
and in L7; \textbf{a second collision} on the area limit law, \S\ref{sec:moments}.
The semiperimeter series, the $60$-term area series and the directed-convex area
series were searched against OEIS on 2026-08-17 and are absent, with a control
query that hits \oeis{A005436} exactly.
\emph{Human verification:} none, as of this draft.}

\begin{abstract}
\noindent
A \emph{convex polyplet} is a king-connected set of cells of $\Z^2$ in which
every row and every column is a single contiguous run. We report what is now
known about counting them, in the order the paper proves it.

The structural result is a squeeze. Let $\mathcal C$ be \emph{any} class with
$\text{staircase} \subseteq \mathcal C \subseteq \text{HV-convex}$. Then
$\lim \mathcal C(n)^{1/n}$ exists and equals the growth constant
$\mu = 3.128943269730886\dots$ of the full class. Reading an HV-convex animal
left to right, it fattens, shears, then thins; the shearing middle is a
staircase polyplet and the two ends are stacks, which are too few to move an
exponential. The proof uses no analytic input, and because the staircase join is
exactly injective the limit is a supremum, so every enumerated term is a rigorous
floor under $\mu$.

By \textbf{area} the class is intractable, and we make that quantitative. On
$700$ terms the generating function satisfies no P-recurrence of order $\le 24$
with polynomial coefficients of degree $\le 24$, and no algebraic equation of
degree $\le 20$ in $F$ with coefficients of degree $\le 20$ in $t$; the same
holds for the classical convex-polyomino control regenerated to the same length.
We pin $\mu$ to $199$ digits and measure the correction to pure exponential
growth to be geometric: $\theta = 0$, no power-law factor and nothing resembling
a stretched exponential. Integer-relation searches find no polynomial with $\mu$
as a root of degree $\le 20$ and height $\le 10^8$, $\le 30$ and $\le 10^6$, or
$\le 12$ and $\le 10^{15}$.

Behind both statistics is one kernel. Solving the area-refined functional
equation by Temperley's method expresses the generating function through a
$q$-Bessel-type function $K$, whose smallest positive zero gives $\mu$ and whose
residue gives the amplitude; exact rational interval arithmetic certifies
$44$ digits of $\mu$ and of that amplitude, agreeing with the extrapolation
above. The same kernel makes the area-moment generating functions algebraic at
every level --- from which the area limit law follows, and that law is already
Richard's.

By \textbf{semiperimeter} the same class is algebraic, of degree $2$; imposing a
directedness condition strong enough to cut the population to about $1\%$ raises
the degree to $4$ and does not leave the algebraic class. Convexity is a
perimeter lever rather than an area lever, for polyplets as for polyominoes.

What sits underneath the leading term is the subject of companion paper L7,
\emph{The subdominant exponential and the amplitude ratio of convex
polyplets}: the four-cone-directed subclass splits into an ascending half growing
at $\mu$ and a descending half growing at a second constant
$\nu = 2.5145796\dots$, and the amplitude ratio of the two classes has a closed
expression.
\end{abstract}

\tableofcontents

\section{Introduction}

Convexity is the classical tractability lever for polyominoes. Convex polygons
by perimeter have an algebraic generating
function~\cite{delestViennot1984}; convex polyominoes by area grow at
$2.30914\dots$~\cite{bender1974}; symmetry classes of convex polyominoes have
been enumerated on the square lattice~\cite{lrr1998} and on the honeycomb
lattice~\cite{gblRoux2004}. It is natural to expect the lever to work on the
king lattice too.

It does, but only for perimeter. The same split appears on the square lattice,
where the classical solvability results are perimeter results and the area side
is not known to be tractable either.

\begin{definition}
A \emph{convex polyplet}, equivalently an \emph{HV-convex polyplet} or convex
animal on the king lattice, is a finite king-connected set of cells of $\Z^2$ in
which every row and every column is a single contiguous run, counted up to
translation. Write $A(n)$ for the number with $n$ cells:
\[
  1,\; 4,\; 16,\; 61,\; 221,\; 766,\; 2566,\; 8390,\; 26982,\; 85834,\;
  271174,\; 853111,\; \dots
\]
A \emph{staircase} polyplet --- the king-lattice analogue of the parallelogram
polyominoes --- is one in which, scanning columns left to right,
both the bottom and the top of the column interval are nondecreasing. Write
$M(n)$ for their number: $1, 3, 9, 28, 87, 272, \dots$ (OEIS \oeis{A225114}, the skew shapes with no empty row or column).
\end{definition}

\section{The class and its transfer structure}
\label{sec:structure}

An HV-convex polyplet is a sequence of column intervals $[b(j), t(j)]$,
$j = 1,\dots,k$, with $b$ valley-unimodal and $t$ peak-unimodal, counted up to
translation. Write
\[
  d(j) = b(j{+}1) - b(j), \qquad h(j) = t(j) - b(j) + 1, \qquad
  s(j) = h(j) - h(j{+}1).
\]

Give column $j$ the \emph{phase} $(p_b, p_t)$, where $p_b = 1$ if some earlier
step had $d > 0$ and $p_t = 1$ if some earlier step had $t$ decreasing.
Unimodality makes both bits monotone in $j$: each flips at most once and never
flips back. So the transfer operator is block-triangular in the four phases, and
the diagonal blocks are:

\begin{table}[H]
\centering\small
\begin{tabular}{c l l l}
\toprule
phase & meaning & step constraint & kernel $K(h,h')$ \\
\midrule
$(0,0)$ & $b$ falls, $t$ rises & $s \le d \le 0$ & $h'-h+1$, \; $h' \ge h$ \\
$(1,0)$ & both rise & $d \ge \max(0,s)$ & $\min(h,h')+1$ \\
$(0,1)$ & both fall & $d \le \min(0,s)$ & $\min(h,h')+1$ \\
$(1,1)$ & $b$ rises, $t$ falls & $0 \le d \le s$ & $h-h'+1$, \; $h' \le h$ \\
\bottomrule
\end{tabular}
\caption{The four diagonal blocks. Blocks $(1,0)$ and $(0,1)$ are the staircase
kernel and its mirror; blocks $(0,0)$ and $(1,1)$ are stacks.}
\label{tab:blocks}
\end{table}

The unbounded state variable is the column height, and the coupling in it is
the single kernel $\min(h,h')+1$, appearing once in each middle block. That
kernel is the whole of the king lattice's difference from the square one, where
it is $\min(h,h')$, and the sections that follow are arguments about it.

\section{The squeeze}
\label{sec:squeeze}

\begin{proposition}[the squeeze]
\label{prop:squeeze}
Let $\mathcal C$ be any class of polyplets with
\[
  \text{staircase} \;\subseteq\; \mathcal C \;\subseteq\; \text{HV-convex}.
\]
Then $\lim_n \mathcal C(n)^{1/n}$ exists and equals
$\lim_n A(n)^{1/n} = \mu = 3.128943269730886\dots$
\end{proposition}

The staircase polyplets are a thin subclass of the HV-convex ones and grow at the
same exponential rate, so every class between them is trapped term by term and
nothing has to be proved about the intermediate classes individually.

\begin{lemma}[three-block factorisation]
\label{lem:blocks}
Write $P(n)$ for the number of monotone-height blocks of area $n$ (phase
$(0,0)$, or equivalently by reversal phase $(1,1)$) and $M(n)$ for the staircase
count, with $P(0) = M(0) = 1$. Then
\begin{equation}
  A(n) \;\le\; 2\,(n{+}1)^2 \sum_{i+j+l=n} P(i)\,M(j)\,P(l) .
  \label{eq:factor}
\end{equation}
\end{lemma}

\begin{proof}
By the monotonicity of the phase bits, the columns split into three consecutive
(possibly empty) runs $C_1C_2C_3$ carrying phases $(0,0)$; one of $(1,0)$ or
$(0,1)$; and $(1,1)$. Every step internal to a run obeys that run's constraint
in Table~\ref{tab:blocks}, so taken standalone $C_1$ and $C_3$ are
monotone-height blocks and $C_2$ is a staircase polyplet or the vertical mirror of
one. The animal is recovered from $(C_1,C_2,C_3)$, the two junction offsets and
one bit naming $C_2$'s phase; each junction offset lies in $[-h', h]$ with
$h + h' \le n$, hence takes at most $n+1$ values.
\end{proof}

\begin{remark}[attribution --- 1 of 3]
\label{rem:attr1}
This decomposition is not new in kind. Gouyou-Beauchamps and
Leroux~\cite{gblRoux2004}, \S2.3 ``Growth phases of convex polyominoes'',
decompose a convex polyomino into blocks $H_{ij}$ by the growth phases of its
upper and lower profiles, with the column state an ordered pair, the transitions
one-way, the extreme blocks $H_{00}$ and $H_{22}$ identified as \emph{stack}
polyominoes and the middle blocks as \emph{staircase} polyominoes with
$H_{02} = Pa = H_{20}$. That is Lemma~\ref{lem:blocks} here, together with
Lemma~\ref{lem:stacks}'s identification, and the mirror equality that companion
paper L7 attributes to them, for convex polyominoes on the honeycomb lattice.
Their square-lattice companion~\cite{lrr1998} reaches the same objects by a
different route, so 2004 is the source to cite. What is not theirs: the king
lattice, where their middle kernel $\min(h,h')$ becomes $\min(h,h')+1$; a class
with no exact solution; and what Proposition~\ref{prop:squeeze} and L7 do with
the blocks.
\end{remark}

\begin{lemma}[the outer blocks are stacks, hence sub-exponential]
\label{lem:stacks}
$P(n)$ is the number of weakly unimodal compositions of $n$ (OEIS
\oeis{A001523}), and $P(n) \le E(n) := (n{+}1)^{4\sqrt n + 6}$, so
$P(n)^{1/n} \to 1$.
\end{lemma}

Only the counting bound is used downstream; the bijection identifies what is
being counted.

\begin{proof}[Part 1: a stack is a weakly unimodal composition]
In a phase-$(1,1)$ block the bottoms rise and the tops fall, so the column
intervals are \emph{nested}: $[b(j{+}1),t(j{+}1)] \subseteq [b(j),t(j)]$.
Nesting means that the set of columns meeting row $y$ is never broken --- it is
a prefix $1..r(y)$. Transpose the picture and record the row widths $r(y)$ from
the bottom up. Since $\{y : r(y) \ge k\} = [b(k),t(k)]$ is a decreasing nested
family of intervals, $r$ rises then falls, and $\sum_y r(y) = n$. The stack is
rebuilt from $r$ alone, so the map is a bijection.

Worked instance: columns $[0,4],[1,3],[2,2]$ (heights $5,3,1$, area $9$) have
row widths $r = (1,2,3,2,1)$, unimodal and summing to $9$.
\end{proof}

\begin{proof}[Part 2: there are few of them]
Cut a weakly unimodal composition at its peak; what falls out on either side is
a partition, so $P(n) \le (n{+}1)^2 p(n)^2$, the $(n{+}1)^2$ paying for the
peak's position and value.

For $p(n)$ a deliberately crude bound suffices, since the squeeze uses only
$P(n)^{1/n}\to1$. Put $s = \lceil\sqrt n\rceil$ and split a partition of $n$ at
$s$. The parts $\le s$ are determined by their $s$ multiplicities, each an
integer in $[0,n]$, contributing at most $(n{+}1)^s$ choices. The parts $> s$
number at most $L = \lfloor n/(s{+}1)\rfloor$, and listing them in nonincreasing
order gives at most $(L{+}1)n^L \le (n{+}1)^{L+1}$ choices. Hence
$p(n) \le (n{+}1)^{s+L+1} \le (n{+}1)^{2\sqrt n + 2}$, since $s \le \sqrt n + 1$
and $L \le \sqrt n$. So $P(n) \le E(n)$ and
$\log P(n)/n \le (4\sqrt n + 6)\log(n{+}1)/n \to 0$.
\end{proof}

Only the exponent matters: $\sqrt n \log n$ is $o(n)$, so the stacks carry no
exponential. Hardy and Ramanujan~\cite{hardyRamanujan1918} give the sharper
$p(n) = e^{O(\sqrt n)}$, which nothing downstream needs, and the elementary
bound keeps the squeeze free of analytic input.

\begin{lemma}[the staircase count is supermultiplicative]
\label{lem:supermul}
$M(i)M(j) \le M(i+j)$ for all $i,j \ge 0$. Hence $\lim_n M(n)^{1/n}$ exists and
equals $\sup_n M(n)^{1/n}$.
\end{lemma}

\begin{proof}
Given staircase polyplets $X$ of area $i$ with last column height $h$, and $Y$ of
area $j$ with first column height $h'$, slide $Y$ up by $d = \max(0, h-h')$ and
butt the two column sequences together.

\emph{Why that $d$ and no other.} Staircase means both boundaries
nondecreasing. At the seam, $b$ not dropping forces $d \ge 0$; $t$ not dropping
forces $d \ge h - h'$. So $\max(0, h-h')$ is the smallest legal slide --- it is a
\emph{function} of $X$ and $Y$, not a choice, which is what stops the join from
smuggling in information.

\emph{The join is in the class.} $d \ge 0$ keeps bottoms nondecreasing;
$d - s = \max(0,h-h') - (h-h') \ge 0$ keeps tops nondecreasing; and $d \le h$,
since $d = h-h' \le h-1$ when $h \ge h'$ and $d = 0$ otherwise, so the junction
columns are king-adjacent. Area is $i+j$.

\emph{The join is injective at fixed $(i,j)$.} Every column is nonempty, so the
cumulative areas $h(1)+\dots+h(r)$ are \emph{strictly} increasing, and a
strictly increasing sequence hits $i$ at most once. Cut at that unique prefix:
those are $X$'s columns and the rest are $Y$'s, each read back up to
translation, which is how both were counted. No split index has to be supplied
alongside the join, so distinct pairs have distinct images.

Fekete's lemma in supermultiplicative form finishes.
\end{proof}

\begin{corollary}[every banked term is a floor]
\label{cor:floor}
$\mu \ge M(n)^{1/n}$ for every $n$, with no extrapolation. At $n = 700$,
\[
  \mu \;\ge\; M(700)^{1/700} \;=\; 3.12340450886853853211\dots
\]
\end{corollary}

The digits past $3.12$ in the extrapolated value remain extrapolation; the
inequality above is exact arithmetic on an enumerated term.

\begin{proof}[Proof of Proposition~\ref{prop:squeeze}]
Write $\mu = \lim M(n)^{1/n} = \sup_n M(n)^{1/n} \ge 1$, which exists by
Lemma~\ref{lem:supermul}. Fix $\varepsilon > 0$ and take $C_\varepsilon$ with
$M(j) \le C_\varepsilon(\mu+\varepsilon)^j$. In \eqref{eq:factor}, bound
$P(i), P(l) \le E(n)$ by Lemma~\ref{lem:stacks}, valid for every $i, l \le n$
since $E$ is nondecreasing:
\[
  A(n) \;\le\; 2(n{+}1)^4 E(n)^2 C_\varepsilon (\mu+\varepsilon)^n ,
\]
and $E(n)^{1/n} \to 1$, so $\limsup A(n)^{1/n} \le \mu + \varepsilon$ for every
$\varepsilon$. In the other direction, staircase $\subseteq \mathcal C
\subseteq$ HV-convex gives $M(n) \le \mathcal C(n) \le A(n)$ termwise, so
$\liminf \mathcal C(n)^{1/n} \ge \mu$ and
$\limsup \mathcal C(n)^{1/n} \le \limsup A(n)^{1/n} \le \mu$.
\end{proof}

\subsection{Consequences of the squeeze}

The four-cone-directed subclass shares $\mu$ with the unrestricted one to $49$
digits, which had looked like a bijection waiting to be found. It is not: it is
Proposition~\ref{prop:squeeze}, and the four-cone condition is not special.
\emph{Every} condition that leaves the staircase polyplets alone leaves $\mu$
alone, because the staircase polyplets already carry all of the exponential
entropy. Table~\ref{tab:blocks} says the same thing at the operator level: the
four-cone floor rewrites the kernels of blocks $(0,0)$ and $(0,1)$ and leaves
blocks $(1,0)$ and $(1,1)$ untouched entry for entry --- and block $(1,0)$ is the
staircase kernel.

\begin{table}[H]
\centering\small
\begin{tabular}{l l r l}
\toprule
class & first terms & trusted digits of $\mu$ & amplitude $C$ \\
\midrule
staircase (\oeis{A225114}) & $1,3,9,28,87,272$ & 204 & $0.28932146397165132308$ \\
HV-convex, $b$ nondecreasing & $1,3,10,33,107,342$ & 202 & $0.37545302027992473174$ \\
(dir4, HV-convex) & $1,4,15,53,177,567$ & --- & $0.45030318571234118235$ \\
HV-convex & $1,4,16,61,221,766$ & 199 & $0.97445221313500464915$ \\
\bottomrule
\end{tabular}
\caption{Four classes, one constant. The staircase value reproduces all $199$
digits the HV-convex series pins, and continues to $204$.}
\label{tab:fourrungs}
\end{table}

\section{Counting by area}
\label{sec:exclusions}

\begin{theorem}
\label{thm:excl}
On $700$ terms, the by-area generating function of HV-convex polyplets
satisfies no P-recurrence
$\sum_{i \le J} p_i(n)a(n{+}i) = 0$ with $J \le 24$ and $\deg p_i \le 24$, and no
algebraic equation $\sum_{j \le K} q_j(t)F(t)^j = 0$ with $K \le 24$ and
$\deg q_j \le 24$. The same holds for HV-convex polyominoes by area
(\oeis{A067675}), regenerated to the same length.
\end{theorem}

Two facts make this one test rather than a $25\times25$ sweep.

\paragraph{The boxes nest.} A solution with $J' \le J$ and $D' \le D$ is a
solution of the $(J,D)$ ansatz with the unused $p_i$ set to zero. Excluding the
maximal box excludes every box inside it.

\paragraph{Full column rank mod $p$ is a proof over $\Q$.} The matrix entries
are integers. Full column rank modulo $p$ means some maximal minor is nonzero
mod $p$, hence nonzero over $\Q$, hence the only rational solution is trivial.
Rank can drop mod $p$, never rise, so \emph{this} direction is rigorous. A rank
\emph{defect} would only be a candidate --- which is the branch that gets the
holdout and the second prime.

\begin{table}[H]
\centering\small
\begin{tabular}{l l l r r r l}
\toprule
series & mode & box & unknowns & rows & rank & verdict \\
\midrule
king    & prec & $(20,20)$ & 441 & 680 & 441 & EXCLUDED \\
king    & prec & $(24,24)$ & 625 & 676 & 625 & EXCLUDED \\
king    & prec & $(20,20)$, skip 50 rows & 441 & 630 & 441 & EXCLUDED \\
king    & prec & $(20,20)$, second prime & 441 & 680 & 441 & EXCLUDED \\
king    & alg  & $(20,20)$ & 441 & 701 & 441 & EXCLUDED \\
king    & alg  & $(24,24)$ & 625 & 701 & 625 & EXCLUDED \\
control & prec & $(20,20)$ & 441 & 680 & 441 & EXCLUDED \\
control & prec & $(24,24)$ & 625 & 676 & 625 & EXCLUDED \\
control & prec & $(20,20)$, skip 50 rows & 441 & 630 & 441 & EXCLUDED \\
control & prec & $(20,20)$, second prime & 441 & 680 & 441 & EXCLUDED \\
control & alg  & $(20,20)$ & 441 & 701 & 441 & EXCLUDED \\
control & alg  & $(24,24)$ & 625 & 701 & 625 & EXCLUDED \\
\bottomrule
\end{tabular}
\caption{$p = 2^{61}-1$ unless noted. The row skip tests ``holds eventually''
rather than ``holds from $n=1$''; same verdict.}
\label{tab:excl}
\end{table}

$(25,25)$ is the first box that $700$ terms cannot decide, at $676$ unknowns
against $675$ rows, so the boxes above are where the term count runs out rather
than where the method does.

Algebraicity was tested separately even though algebraic implies D-finite, since
the implication only rules out algebraic functions whose induced ODE fits inside
the box already excluded, and a degree-$24$ algebraic function can induce a much
larger one.

\subsection{Positive and negative controls}

A negative from a guesser means nothing without a positive control: a series
whose answer is already known, run through the identical pipeline, where the
tool has to return that answer. Four arms run as a gate, each required to come
out as stated:

\begin{table}[H]
\centering\small
\begin{tabular}{l l l l l}
\toprule
arm & series & box & must be & is \\
\midrule
prec $+$ & HV-convex by semiperimeter & $(5,2)$ & CANDIDATE & nullity 1; predicts all 172 held-out rows \\
prec $-$ & same series & $(2,1)$ & EXCLUDED & rank 6 of 6 \\
alg $+$ & \oeis{A005436}, algebraic~\cite{delestViennot1984} & $(2,8)$ & CANDIDATE & nullity 1; holds on all 69 held-out rows \\
alg $-$ & same series & $(2,4)$ & EXCLUDED & rank 15 of 15 \\
\bottomrule
\end{tabular}
\caption{The tool finds a P-recurrence when one exists, finds an algebraic
equation when one exists, and refuses both when the box is too small.}
\label{tab:controls}
\end{table}

A fifth arm is what makes any positive verdict in this paper trustworthy: a
series with no algebraic equation to find. $199$ terms of the by-area king
series, the same length and comparable term size to the perimeter series, come
back EXCLUDED with nullity $0$ at every box where the perimeter series is a
candidate. A $199$-term matrix of that shape does not manufacture rank
defects.

\subsection{$\mu$, $\theta$, and the amplitude}
\label{sec:mu}

The discriminator for the shape of the correction is the ratio of successive
differences of $r_n = a(n{+}1)/a(n)$: a constant below $1$ means the correction
is geometric --- an isolated subdominant singularity, no power-law factor ---
while drift toward $1$ would mean a branch point and $\theta \ne 0$.

\begin{table}[H]
\centering\small
\begin{tabular}{l l l}
\toprule
series & $d_n/d_{n-1}$ & reading \\
\midrule
king (HV-convex) & $0.48100879371$, constant to all 12 printed digits & geometric \\
control (\oeis{A067675}) & $0.625135968591$, constant to all 12 printed digits & geometric \\
\bottomrule
\end{tabular}
\caption{Measured over the last five $n$ of each $700$-term series.}
\label{tab:theta}
\end{table}

The direct probe agrees: $n(r_n/\mu - 1)$, which tends to $\theta$ under a power
law, sits at the $10^{-216}$ level for the king series and $10^{-137}$ for the
control --- that is, at zero. So $\theta = 0$ for both, there is no power-law
correction, and nothing resembling a stretched exponential:
\[
  a(n) = C\mu^n\bigl(1 + O(\rho^n)\bigr), \qquad
  \rho = 0.481008794 \text{ (king)}, \quad 0.625135969 \text{ (control)}.
\]

Trusted digits are \emph{measured}, not asserted: the whole pipeline is re-run
on the series truncated by $30$ terms, and the reported figure is the minimum
over three independent comparisons, minus a two-digit guard.

\begin{quote}\small
$\mu$, convex polyplets by area, $199$ trusted digits:\\[2pt]
\texttt{3.1289432697308862522774479953877541605320912219043941349649746499492}\\
\texttt{4438583718257615823063288164783234634835221018937039608167576493048146}\\
\texttt{233778414064847543298805623983850111109827008627878918922549}
\end{quote}

\begin{quote}\small
$\mu$, convex polyominoes by area (\oeis{A067675}), $121$ trusted digits:\\[2pt]
\texttt{2.3091385933304947310987203050172125319118144725816284016944029002844}\\
\texttt{56440748316842717281615774412174374610237798122137142}
\end{quote}

The control's value is an external check on the whole apparatus: Bender's
published constant is $2.30914$~\cite{bender1974}, and the non-king transfer
matrix plus this extrapolation reproduce it and extend it by $116$ digits.

\subsection{The kernel, and what it certifies}
\label{sec:kernel}

The measurements above extrapolate from terms. There is also a closed route to
the same constants, and it certifies its own digits.

Solving the area-refined functional equation by Temperley's method --- the
$q$-deformation of the semiperimeter solution of \S\ref{sec:perimeter} ---
writes the full generating function as an explicit finite combination of four
$q$-adically convergent $q$-series, from which coefficient extraction is exact.
Its denominator telescopes to a single $q$-Bessel-type function,
\[
  K(q) \;=\; \sum_{m \ge 0} (-1)^m\,(2 - q^m)\,
             \frac{q^{m(m+1)/2}}{(q;q)_m^{2}},
  \qquad (q;q)_m = \prod_{j=1}^{m}(1 - q^{j}),
\]
and if $q_c$ is its smallest positive zero then $\mu = 1/q_c$, with the
amplitude given by a residue formula in $K'(q_c)$ and the numerator series.
Dropping the two king boundary terms gives the convex-polyomino control, whose
kernel is the Klarner--Rivest function; that pipeline reproduces \oeis{A067675}
and \oeis{A067676} to $50$ terms each.

\begin{quote}\small
$q_c = 0.319596718059387465518602891982713923614508377\dots$, certified to $45$
digits; $\mu = 3.12894326973088625227744799538775416053209122\dots$, certified
to $44$; $A = 0.9744522131350046491513294208602433274253844580\dots$, certified
at the same precision with a simple pole at $q_c$ and a nonzero residue.
\end{quote}

Certified means exact rational interval arithmetic with no floating point in
the chain: alternating-series bracketing of $K$, an analytic positivity bound on
$(0, 29/100]$, a certified Lipschitz bound on $[0.29, 0.33]$, and a $270$-step
march that pins $q_c$ inside an interval of width $4\times10^{-46}$ and
establishes that it is the smallest positive zero. Those $44$ digits agree with
every corresponding digit of the $199$ extrapolated above, which is the strongest
cross-check in this paper: two methods sharing no machinery --- one series
extrapolation, one interval-arithmetic certificate --- and no disagreement.

The amplitude is the part that changes status. \S\ref{sec:mu} measures $A$; the
residue formula \emph{derives} it, and the derivation makes $A$ a certified
number rather than a fitted one. What it does not yet do is prove
$a(n) \sim A\mu^n$: that needs the analytic facts about $K$ collected in
Open Problem~\ref{op:kzeros}, not just the value of one residue.

\subsection{Area moments, and where they were already known}
\label{sec:moments}

The same kernel, differentiated, gives the area moments. For every $r \ge 0$ the
bivariate $r$-th area-moment generating function of convex polyplets with
bounding box exactly $w \times h$ is algebraic over $\Q(x,y)$ --- explicitly an
element of $\Q(x,y,\sqrt{\Delta})$ with $\Delta = (1-x-y)^2 - 4xy$, built by a
recursive kernel construction --- so the univariate moment generating functions
by semiperimeter are algebraic at every level.

\textbf{The limit law that follows is published, and was rederived here in
ignorance of that.} Normalised moments give
$\mathrm{area}/s^2 \to U(1-U)/2$, the rectangles area law
$\beta_{1,1/2}$; that is Richard~\cite{richard2007}, Table~1 and the remark
following it, for convex polygons, and Enting and
Guttmann~\cite{entingGuttmann1989} cover the same ground for the square-lattice
control. Neither is cited here for support --- they are cited because they got
there first.

What survives the collision is narrow and stated as such: algebraicity at
\emph{every} moment level via one kernel construction rather than level by
level, and the explicit statement that the king and square lattices carry the
identical limit law despite neither class being solved on the king side.

\subsection{Integer-relation searches for $\mu$}

An integer-relation search on $(1,\mu,\dots,\mu^d)$, run at the trusted
precision only. The discipline that matters: a hit is reported as
REJECTED-ARTIFACT unless $(d{+}1)\log_{10}(\text{height})$ is well below half
the trusted digit count, because PSLQ always returns \emph{something} once the
unknowns can absorb the input precision.

\begin{table}[H]
\centering\small
\begin{tabular}{l r r l l}
\toprule
$\mu$ & digits fed & degree $\le$ & height $\le$ & result \\
\midrule
king & 199 & 20 & $10^{8}$ & none at any degree \\
king & 199 & 30 & $10^{6}$ & none at any degree \\
king & 199 & 12 & $10^{15}$ & none at any degree \\
king & 137 & 20 & $10^{5}$ & none at any degree \\
control & 121 & 12 & $10^{8}$ & none at any degree \\
control & 121 & 20 & $10^{4}$ & none at any degree \\
control & 121 & 20 & $10^{8}$ & hits at degree 15--20, \textbf{all REJECTED-ARTIFACT} \\
\bottomrule
\end{tabular}
\caption{The last row is the box outgrowing the precision: capacity
$125$--$142$ against $121$ trusted digits, which is what the capacity test
rejects.}
\label{tab:pslq}
\end{table}

So the falsifiable statement is: \emph{no integer polynomial has $\mu_{\text{king}}$
as a root with degree $\le 20$ and coefficients below $10^8$, nor degree $\le 30$
below $10^6$, nor degree $\le 12$ below $10^{15}$}. A closed form, if one exists,
lies outside those boxes.

\subsubsection*{A worked rejection}

The cubic $3219x^3 - 7630x^2 - 18493x + 33955$ was offered mid-investigation,
produced by feeding a then-published $16$-digit $\mu$ to a computer algebra
service. Its nearest root is
\[
  \begin{aligned}
    &3.12894326973088625330312785623831607097 \\
    \mu = {}&3.12894326973088625227744799538775416053
  \end{aligned}
\]
--- agreement to $18$ digits, disagreement at the $19$th. It is also inside the
PSLQ box already searched (degree $3$, height $3.4\times10^4$), which is why that
search returned nothing at degree $3$.

\section{Counting by semiperimeter}
\label{sec:perimeter}

By \emph{semiperimeter}
--- for a convex animal the box's $W + H$, with edge-perimeter $2(W{+}H)$ ---
the same class is not merely D-finite but algebraic.

\begin{table}[H]
\centering\small
\begin{tabular}{l l l l}
\toprule
class & statistic & verdict & minimal box \\
\midrule
HV-convex polyplets & area & non-D-finite, non-algebraic & --- \\
HV-convex polyplets & semiperimeter & \textbf{algebraic, degree 2} & $(2,9)$ \\
dir4 HV-convex polyplets & area & non-D-finite, non-algebraic & --- \\
dir4 HV-convex polyplets & semiperimeter & \textbf{algebraic, degree 4} & $(4,22)$ \\
\bottomrule
\end{tabular}
\caption{The split between the two statistics survives a directedness
constraint strong enough to cut the population to about $1\%$ by $s = 200$.}
\label{tab:perim}
\end{table}

The unrestricted quadratic, fitted on $34$ rows and holding on all $166$ later
ones with zero failures over $\Z$:
\[
\begin{aligned}
q_0 &= 2t^2 - 21t^3 + 88t^4 - 196t^5 + 242t^6 - 119t^7 + 72t^8 + 16t^9, \\
q_1 &= -4t + 52t^2 - 252t^3 + 554t^4 - 520t^5 + 96t^6 + 128t^7, \\
q_2 &= 2 - 31t + 176t^2 - 416t^3 + 256t^4 + 256t^5 .
\end{aligned}
\]
$(2,8)$ is EXCLUDED and rational is excluded to $t$-degree $98$, so the box is
minimal in both directions.

The directed quartic has $\deg q_j \le 22$; it was fitted on the first $119$
rows, annihilates all $81$ later ones, and re-substituted into the integer
series annihilates all $200$ rows over $\Z$ with zero failures. $(4,21)$ is
EXCLUDED, as is every algebraic degree below $4$ as far as $199$ terms reach
(degree $3$ to $t$-degree $48$, degree $2$ to $65$, rational to $98$).

Two of the directed subclass's series turn out to be known sequences wearing
other clothes. By semiperimeter it is \oeis{A014300}, the count of nodes of odd
outdegree in ordered rooted trees, shifted by one; the diagonal $d(n,n)$ of its
box refinement is \oeis{A112029} $= \sum_k \binom{n-1+k}{k}^2$, likewise
shifted. Both were checked term by term against the entries' own b-files.
Neither entry carries any polyomino, convex or king interpretation, so what is
new here is the interpretation and not the sequence --- material for a comment
on those entries, not a claim in this paper.

\begin{remark}[the coefficients are themselves evidence]
The quartic's coefficients have five digits. A rank defect absorbed from $199$
large terms would produce sixty-digit integers --- and does: the $(19,3)$
P-recurrence for the same object, which is a \emph{non-minimal} description, has
coefficients of about sixty digits.
\end{remark}

\subsection{The nullity check}

The internal consistency check for a positive algebraicity verdict. A genuine
minimal $(K_0,L_0)$ relation forces nullity $(K-K_0+1)(L-L_0+1)$ as the box
opens, since $F^it^jP$ all lie in the larger box. Measured against predicted for
$(K_0,L_0) = (4,22)$, the two grids agree cell for cell over the whole decidable
region; on the by-area series of the fifth arm the same grid is all zeros.

\subsection{Column-convex polyplets}

For contrast, the neighbouring class is exactly solvable.
Temperley's method~\cite{temperley1956} with the last column's height as
catalytic variable --- king contact gives $h + h' + 1$ placements --- closes to a
$2\times2$ linear system and yields
\[
  F(x) \;=\; \frac{x(1-x)^3}{1 - 7x + 13x^2 - 10x^3 + 2x^4},
\]
with growth $4.64468\dots$, a quartic root. Brute-validated to $n \le 8$.

This is a rediscovery: the generating function is that of OEIS \oeis{A187077}
verbatim. Two things about that entry are recorded here, since it carries no
derivation. First, its comment ``Equivalent to a sequence of row-convex
polyhexes (\oeis{A059716})'' is not correct in its plain reading: row-convex
polyhexes are $1,3,11,42,162,\dots$ while this sequence is $1,4,18,83,385,\dots$,
so the two agree only at the first term. Checked by direct enumeration on both
lattices with all fixed polyhexes reproducing \oeis{A001207} as a control. One
placement is the whole difference --- a row of length $h$ admits $h + h' + 1$
positions for a length-$h'$ row beside it on the king lattice against $h + h'$
on the hexagonal one. Second, what the sequence counts is row-convex
\emph{polyplets}; column-convex polyplets are the transpose and give the same
counts.

\section{Relation to existing work}
\label{sec:related}

\paragraph{Attribution --- 2 of 2.} The block decomposition of
Lemma~\ref{lem:blocks}, and the identification of the extreme blocks as stacks
and the middle blocks as staircases, are present in Gouyou-Beauchamps and
Leroux~\cite{gblRoux2004}, \S2.3, for convex polyominoes on the honeycomb
lattice. Their square-lattice companion~\cite{lrr1998} reaches the same objects
by a different route, through partitions, stacks, shifted stacks and directed
convex via Temperley and Bousquet-M\'elou, so~\cite{gblRoux2004} is the source
to cite for the decomposition itself. The third placement that
\repo{docs/publication-split.md} requires sits with the mirror equality, in
companion paper L7.

What is not theirs, and is what this paper claims: the king lattice, where the
middle kernel is $\min(h,h')+1$ rather than $\min(h,h')$ and neither class is
exactly solved; and the squeeze over \emph{every} intermediate class
(Proposition~\ref{prop:squeeze}).

\paragraph{The square-lattice analogue.} Bender's $2.30914$~\cite{bender1974} for
convex polyominoes by area, and the parallelogram subclass \oeis{A006958}
sharing it (measured here at $2.309138593330495$, flat from $n = 100$ to
$n = 400$), mean that Proposition~\ref{prop:squeeze}'s \emph{conclusion} could
have been read off the solved models one lattice over. It is the generality ---
every intermediate class, on a lattice where nothing is solved --- that is the
contribution.

\paragraph{Perimeter.} Delest and Viennot~\cite{delestViennot1984} is the
classical algebraicity result and is used here as the positive control for the
algebraicity guesser. \S\ref{sec:perimeter}'s finding is the king analogue with
an explicit minimal box.

\paragraph{Area moments.} Richard~\cite{richard2007} and Enting and
Guttmann~\cite{entingGuttmann1989} hold the area limit law for convex polygons,
and \S\ref{sec:moments} says plainly that the derivation here reached it second.
It is recorded because a reader who finds the collision unaided has reason to
distrust everything around it; a claim of priority that a literature pass would
have removed is worse than the omission it papers over.

\paragraph{Column-convex.} \S\ref{sec:perimeter} records a rediscovery of
\oeis{A187077} by Temperley's method~\cite{temperley1956}, together with a
correction to that entry's comment.

\section{Open problems}

\begin{openproblem}[a sharp asymptotic, for anything in this family]
Prove $A(n) \sim C\mu^n$, or the same for $M(n)$. Proposition~\ref{prop:squeeze}
and Lemma~\ref{lem:supermul} give exponential rates and never an amplitude, and
the amplitudes of Table~\ref{tab:fourrungs} are measured, not derived.
\S\ref{sec:kernel} narrows what is missing rather than closing it: the constant
$A$ now comes from a residue at a certified simple pole, so the gap is no longer
the amplitude's value but the analytic bookkeeping around it --- that the
remaining singularities of the kernel quotient are strictly farther out and that
the tail they contribute is $O(\rho^n)$. A second route, independent of the
kernel, is a proof of Kurkov's conjectured continued fraction for
\oeis{A225114}; companion paper L7 states what would follow from it.
\end{openproblem}

\begin{openproblem}[the zeros of $K$, and a route to non-D-finiteness]
\label{op:kzeros}
$K$ appears to have infinitely many real zeros in $(0,1)$ accumulating at
$q = 1$: at least $40$ are located, and near the accumulation point they follow
\[
  K(e^{-\varepsilon}) \sim 2\cdot 3^{1/4}\sqrt{\varepsilon/2\pi}\,
  \cos\!\left(\frac{V}{\varepsilon} - \frac{\pi}{12}\right),
  \qquad V = 2\,\mathrm{Cl}_2(\pi/3),
\]
with $V$ the Gieseking constant. The generating function has poles at these
zeros, with the residue nonzero at the first four checked. A D-finite function
has finitely many singularities, so infinitely many poles would settle
Theorem~\ref{thm:excl} as a theorem rather than an exclusion on $700$ terms.
Two things are missing and neither is bookkeeping: a rigorous
steepest-descent treatment of the oscillation law as $q \to 1^-$, and
nonvanishing of the residue at infinitely many zeros. Firm numerically, and a
lead rather than a result.
\end{openproblem}

\begin{openproblem}[non-D-finiteness as a theorem]
Theorem~\ref{thm:excl} is an exclusion on $700$ terms, not a proof. The
arithmetic obstruction that works for the height-anisotropic king generating
function does not obviously apply here, since the natural slicing of a convex
class is not by height. Open Problem~\ref{op:kzeros} is the more promising
route.
\end{openproblem}

\section{Reproduction}
\label{sec:repro}

\begin{table}[H]
\centering\small
\begin{tabular}{l l}
\toprule
claim & check \\
\midrule
$A(n)$ to $n = 700$ & \texttt{make build/convex\_area\_tm \&\& build/convex\_area\_tm 700 1} \\
the control to $n = 700$ & \texttt{build/convex\_area\_tm 700 0} \\
Theorem~\ref{thm:excl} & \texttt{make build/prec\_guess}; \texttt{build/prec\_guess prec ... 24 24} \\
Table~\ref{tab:controls} & \texttt{make gate-convex-dfinite} \\
$\mu$, $\theta$, the amplitudes & \repo{experiments/convex\_growth.py} \\
Table~\ref{tab:pslq} & \repo{experiments/convex\_growth.py} \texttt{--algdeg} \\
Lemma~\ref{lem:supermul}, measured & \repo{experiments/staircase\_supermul.py} \\
Lemma~\ref{lem:supermul}, in Lean & \texttt{cd polyplets \&\& lake build} \\
Lemma~\ref{lem:blocks}, brute force & \repo{experiments/monotone\_block\_growth.py} \\
the block series $T$ & \repo{experiments/dir4\_descent\_block.py}, \repo{experiments/descent\_block\_oracle.py} \\
\S\ref{sec:perimeter} & \repo{experiments/dir4\_perim\_find\_alg.py}, \texttt{make gate-dir4-perim-alg} \\
column-convex & \repo{experiments/colconvex\_king.py} \\
the squeeze, pinned & \texttt{make gate-middle-kingdom} \\
\bottomrule
\end{tabular}
\caption{Where each claim is checked. Every proof-relevant item is either a
paper proof above or a Lean theorem; everything in this table is a measurement
or a reproduction of one.}
\label{tab:repro}
\end{table}

\bibliographystyle{plain}
\bibliography{shared/refs}

\end{document}
